\begin{abstract}
多个同步单天线接收站对未知宽带波形进行直接定位时，频谱重叠与功率不平衡使位置搜索和信号估计相互耦合。单源空间谱的高分位置未必构成有效的多源解，逐源更新也可能因反复估计信号系数而缓慢收敛。本文在已知自由空间相对幅度与基带时延模型下，以最小二乘消去未知波形，建立统一的位置目标，并提出两种求解策略。多候选联合精修AP保留条件谱中的少量局部极大值，对各分支联合优化位置，再按拟合残差择优；波形消元块NOMP在残差发现和循环更新中使用消元后的局部位置目标，减轻固定信号系数对位置步长的约束。通过推导条件投影与消元曲率，本文解释两种策略的作用，并以配对实验检验其精度与计算成本。已完成的120条双源观测中，多候选联合AP恢复了单候选遗漏的10条观测；在相同停止精度下，单候选联合精修的平均求解时间比交替精修减少约8.2\%。八条观测的受控机制实验进一步表明，消元更新能显著减少固定系数更新的重复计算。候选保留与波形消元分别从搜索可靠性和更新效率改善了该模型下的直接定位。
\end{abstract}
\noindent\textbf{关键词：}直接位置估计；强弱重叠源；交替投影；块NOMP；可分离非线性最小二乘
